\chapter{場面別・設定の早見表}\label{app:cheatsheet}

迷ったときの出発点です。まず既定値で処理し、結果を見てから下の表の方向へ1つずつ動かしてください。

\begin{center}\small
\renewcommand{\arraystretch}{1.25}
\begin{tabularx}{\linewidth}{@{}>{\raggedright\arraybackslash}p{0.27\linewidth} X@{}}
\toprule
こんなとき & 変える設定（→ 方向）\\
\midrule
\multicolumn{2}{@{}l}{\sffamily\bfseries 撮影の条件}\\
ゆらぎが大きい日 & 採用率 ↓（5〜10\%）、位置合わせ領域の大きさ ↓、探索範囲 ↑（±24〜32）\\
ゆらぎが小さい日 & 採用率 ↑（15〜30\%）、探索範囲 ↓（±8〜12）\\
フレームが少ない（数百枚以下） & 採用率 ↑（20〜50\%）、参照フレーム ↑（25〜50\%）\\
暗い・ノイズが多い & 採用率 ↑、ウェーブレットのノイズ ↑、強調 ↓\\
途中で雲・ピント直し & 品質グラフの時系列を見て、使うフレームの範囲で悪い区間を外す\\
\midrule
\multicolumn{2}{@{}l}{\sffamily\bfseries 対象}\\
木星・土星・火星 & 既定値から。ウェーブレット Layer 1〜3 を上げる。RGB チャンネル合わせは \ui{自動}\\
月面・太陽面 & 対象モード \uil{自動}（または\uil{月・太陽}）。ウェーブレットは Layer 1〜2 控えめ、3〜4 を少し\\
カメラの RAW & 仕上げの色 \ui{自動（灰色仮説）} とレベル補正を必ず\\
大きなセンサーで対象が小さい & 処理範囲で対象の周りだけを処理\\
\midrule
\multicolumn{2}{@{}l}{\sffamily\bfseries 結果の見え方}\\
ぼやける & 採用率 ↓、ウェーブレット強調 ↑、位置合わせ領域 ↓\\
ざらつく & 採用率 ↑、ウェーブレットのノイズ ↑\\
ゆがむ・崩れる & 一致度の下限 ↑（0.6〜0.7）、位置合わせ領域 ↑\\
縁に黒い輪 & 輪抑制 ↑（0.2〜0.5）、Layer 1〜2 の強調 ↓\\
縁に赤・青のにじみ & RGB チャンネル合わせ \ui{自動で合わせる}\\
光跡・白い点 & 加算方式 \uil{σクリップ}（閾値 2.0〜3.0）\\
\bottomrule
\end{tabularx}
\end{center}

\chapter{設定の一覧}\label{app:settings}

\begin{center}\footnotesize
\renewcommand{\arraystretch}{1.15}
\begin{tabularx}{\linewidth}{@{}l>{\raggedright\arraybackslash}p{0.27\linewidth}>{\raggedright\arraybackslash}X@{}}
\toprule
タブ・欄 & 設定 & 既定値（範囲）\\
\midrule
品質評価 & 品質指標 & 勾配エネルギー（周波数帯パワー比）\\
 & バイトオーダー（SERのみ） & 自動判定（little／big）\\
 & ビット深度の解釈 & ヘッダに従う（12bit／14bit）\\
 & 処理範囲 & 全体\\
 & 使うフレームの範囲 & 空欄＝全体\\
 & 色の配列（Bayer） & ヘッダに従う（モノクロ／RGGB／GRBG／GBRG／BGGR）\\
 & デバイヤーの方式 & 標準 bilinear（高品質 Malvar-He-Cutler）\\
 & ダーク・フラット & なし\\
 & 外れ値判定の厳しさ & 6.0 σ\\
 & 参照との類似度の下限 & 0.5\\
 & 許容する位置ずれ & 自動（短い辺の 1/4）\\
\midrule
アライメント & 位置合わせ方法 & 複数領域の局所位置合わせ（画像全体のみ）\\
 & 対象モード & 自動（惑星／月・太陽）\\
 & 位置合わせ領域の大きさ & 自動（32・48・64・96・128・200 px）\\
 & 位置ずれの探索範囲 & ±16 px（±8〜±32）\\
 & 参照フレーム & 品質上位 25\%（5〜50\%）\\
 & 参照の反復精密化 & オン\\
 & 一致度の下限 & 0.5\\
 & 配置する模様の強さ・明るさ & 0.6・0.15\\
\midrule
スタック & フレームの選択方式 & 割合（枚数）\\
 & 採用するフレーム & 10\%（1〜100\%）\\
 & 輝度正規化 & オン\\
 & 加算方式 & 単純平均（品質重み付き平均／σクリップ）\\
 & σクリップの閾値 & 2.0\\
 & 低メモリモード & オフ\\
 & Bayer のまま合成 & オフ\\
 & ドリズルの倍率 & 1×（1.5×／2×／3×）\\
 & pixfrac & 0.90（0.5〜1.0）\\
\midrule
仕上げ・出力 & RGB チャンネル合わせ & 0（画素、小数可）\\
 & ウェーブレットの強調 & 1.00（0〜20）\\
 & ウェーブレットのノイズ & 0.00（0〜1）\\
 & 連動の強さ・輪抑制 & 1.00（0〜2.5）・0.00（0〜1）\\
 & R・G・B の倍率・彩度 & 1.0（0.5〜2.0）・1.0（0〜2）\\
 & レベル補正 & 黒 0・中間 1.00・白 255\\
 & 向き・切り抜き & なし\\
 & 形式 & 16bit TIFF（32bit float TIFF／32bit float FITS／16bit PNG）\\
 & ファイル名の付け方 & 元の名前＋処理条件（WinJUPOS形式）\\
 & 処理条件と撮影時刻の記録 & オン\\
 & 複数の採用率 & 5, 10, 25\\
\bottomrule
\end{tabularx}
\end{center}

\chapter{キーボードとマウスの操作}\label{app:keys}

\begin{center}\small
\renewcommand{\arraystretch}{1.2}
\begin{tabularx}{\linewidth}{@{}>{\raggedright\arraybackslash}p{0.33\linewidth} X@{}}
\toprule
操作 & はたらき\\
\midrule
\key{⌘}\key{O} & ファイルを開く（入力キューへ追加）\\
\key{⌘}\key{E} & 品質評価\\
\key{shift}\key{⌘}\key{A} & アライメント\\
\key{⌘}\key{R}（または \key{return}） & スタック\\
\key{⌘}\key{.} & 処理を中断\\
\key{⌘}\key{S} & 書き出し\\
\key{shift}\key{⌘}\key{S} & 複数の採用率で書き出し\\
\key{⌘}\key{1}・\key{⌘}\key{2} & 左の欄・右の設定パネルを隠す／出す\\
\key{⌘}\key{+}・\key{⌘}\key{-} & 拡大・縮小\\
\key{⌘}\key{0}・\key{⌘}\key{9} & 全体を表示・実画素で等倍（100\%）\\
\key{⌘}\key{Z}・\key{shift}\key{⌘}\key{Z} & 位置合わせ領域の編集などの取り消し・やり直し\\
\midrule
\key{←}\key{→} & フレームを1枚送る（\key{shift} で10枚、\key{option} で100枚）\\
\key{⌘}＋ホイール、ピンチ & カーソルの位置を中心に拡大・縮小\\
2本指のスクロール・ホイール & 拡大しているときに見る場所を動かす\\
品質グラフをクリック & そのフレームを表示\\
品質グラフのオレンジの線をドラッグ & 参照に使う割合を変える\\
（配置を編集のとき）クリック・\key{delete} & 位置合わせ領域を足す・選んで消す\\
入力キューを右クリック & Finder で表示・キューから削除\\
\bottomrule
\end{tabularx}
\end{center}

\chapter{この説明書について}

この説明書の画面は、\LS \version\ を実際に動かして撮ったものです。画面の説明の番号や枠の位置は、アプリが書き出した部品の位置から自動で置いています。
アプリの画面が変わったときは、\texttt{docs/manual/tools/make\_shots.sh} で撮り直し、\texttt{tools/figures.py} で図を作り直してから組版します
（手順は \texttt{docs/manual/README.md}）。
